\section{Conceptual Architecture}
\label{sec:conceptual}

\sectionnote{Owner: Jeremy McCrindle. Target: 1.25 pages plus the box-and-arrow diagram. Rubric: Conceptual architecture breakdown is clearly referenced, including rationale for updating the architecture (compared to A1); Good subsystem layouts.}

\subsection{Starting Point: The Conceptual Architecture Report}
\label{sec:a1-summary}

Our conceptual architecture report~\cite{a1-report} described WezTerm through eight subsystems, identified by tracing SSH multiplexing use cases. The GUI Front End, Multiplexer, Mux Client and SSH Client ran on the local host, and the Mux Server, Terminal Emulation Core and PTY ran on the remote host. A Configuration Engine was described in the text but not drawn. The report said little about fonts, windowing or scripting, which every session uses, so this report extends the architecture to the whole system.

\subsection{Changes Since the Conceptual Architecture Report}
\label{sec:a1-changes}

Each change is based on WezTerm's documentation, the conceptual architecture reports of other groups, or inconsistencies in our own report. None is based on the source code, because differences from the code are the divergences analysed in Sections~\ref{sec:reflexion} and~\ref{sec:subsystem}. From group 22~\cite{g22-a1} we adopted font handling as a separate component and configuration as a component that scripts can call, and from group 13~\cite{g13-a1}, domains, which make local and remote panes behave alike. Table~\ref{tab:components} lists the change to each component; none was removed.

\begin{figure}[!bt]
  \centering
  \begin{minipage}{0.85\linewidth}\diagram{conceptual-architecture.tex}\end{minipage}
  \caption{Revised conceptual architecture of WezTerm. An arrow points from a component to the component it depends on, and dashed arrows are configuration reads. Below the client-server line, the server runs its own Multiplexer, Terminal Emulation Core and PTY, which are drawn once. For readability, the figure omits User $\rightarrow$ Mux Client (\texttt{wezterm cli}), GUI Front End $\rightarrow$ Mux Server and three configuration reads, all of which are listed in Table~\ref{tab:dependencies}.}
  \label{fig:conceptual}
\end{figure}

The main dependency change concerns the PTY. Our earlier report showed the Terminal Emulation Core depending on the PTY~\cite{a1-report}. The Multiplexer now depends on it instead, since the domain that creates a pane also creates its PTY~\cite{what-is-a-terminal,multiplexing}, and all three other groups model it the same way~\cite{g7-a1,g13-a1,g22-a1}. The Terminal Emulation Core only reads and writes bytes, so the earlier arrow showed a flow of data, not a dependency. For the same reason, the Mux Server now depends on the Multiplexer instead of the Terminal Emulation Core.

The revision also fixes two inconsistencies in our earlier report. The Configuration Engine was described but not drawn, and it is now drawn. The local-pane use case placed a Mux Server and Terminal Emulation Core on the local host, while the figure showed both only on the remote host. The revised figure has no host boundaries, since the same components serve local panes and run inside a remote Mux Server.

\subsection{Revised Conceptual Architecture}
\label{sec:revised-conceptual}

Figure~\ref{fig:conceptual} shows the 11 components and five external entities, with a local command following the centre column and the remote path on the right. The architecture keeps the two styles of our earlier report and adds a third~\cite{arch-styles}. In the client-server style, the Mux Client and SSH Client are clients of the Mux Server and SSH Server, and the Mux Protocol defines how the Mux Client and Mux Server communicate~\cite{multiplexing}. In the implicit invocation style, the GUI Front End and Multiplexer announce events, such as new output in a pane, to listeners they do not know~\cite{event-registration}, so the figure has no arrow from the Multiplexer to the GUI Front End. The new interpreter style covers the Configuration and Lua Engine, an embedded Lua 5.4 interpreter that runs the user's \texttt{wezterm.lua}~\cite{lua-reference}.

Table~\ref{tab:components} compares each component with the web-browser reference architecture~\cite{refarch-browser}, the closest published one. Like a browser, a terminal parses marked-up content into a model, displays it and can be extended by scripts. WezTerm differs in reaching its content sources through the Multiplexer, the counterpart of the browser engine, and in having a server side. Table~\ref{tab:dependencies} explains each dependency. Dependencies 14 and 19 are inferred from the documentation.

\begin{table}[!htbp]
  \centering
  \caption{Components of the revised conceptual architecture, with each component's change from the conceptual architecture report.}
  \label{tab:components}
  \footnotesize
  \setlength{\tabcolsep}{4pt}
  \begin{tabularx}{\linewidth}{@{}>{\raggedright\arraybackslash}p{0.215\linewidth}>{\raggedright\arraybackslash}X>{\raggedright\arraybackslash}p{0.23\linewidth}>{\raggedright\arraybackslash}p{0.15\linewidth}@{}}
    \toprule
    \textbf{Component} & \textbf{Role} & \textbf{Change} & \textbf{Browser counterpart} \\
    \midrule
    GUI Front End & Handles input and draws panes & Narrowed & User Interface \\
    Window and Platform Layer & Provides native windows and events & New; was inside the GUI & Display Backend \\
    Font Subsystem & Finds, shapes and rasterises fonts & New; split from the GUI & Display Backend \\
    Multiplexer & Holds the sessions and creates panes & Now owns local panes & Browser Engine \\
    Terminal Emulation Core & Parses each pane's bytes into a cell grid & Runs with every Multiplexer & Rendering Engine \\
    PTY & Runs each program in a pseudo-terminal & Now also used locally & Networking \\
    Configuration and Lua Engine & Runs \texttt{wezterm.lua} and its handlers & Renamed and now drawn & JavaScript Interpreter \\
    Mux Client & Displays a Mux Server's panes locally & Widened beyond SSH & Networking \\
    Mux Protocol & Defines the client-server messages & New; formerly implicit & Networking \\
    Mux Server & Serves a Multiplexer's panes to clients & Serves its own Multiplexer & Web server \\
    SSH Client & Built-in SSH client (libssh or libssh2) & Adds plain SSH domains & Networking \\
    \bottomrule
  \end{tabularx}
\end{table}

\begin{table}[!htbp]
  \centering
  \caption{Dependencies of the revised conceptual architecture. Dependencies 15, 16 and 19--23 are configuration reads.}
  \label{tab:dependencies}
  \footnotesize
  \begin{tabularx}{\linewidth}{@{}l>{\raggedright\arraybackslash}p{0.53\linewidth}>{\raggedright\arraybackslash}X@{}}
    \toprule
    \textbf{\#} & \textbf{Dependency} & \textbf{Reason} \\
    \midrule
    1 & GUI Front End $\rightarrow$ Multiplexer & The GUI displays and controls the session \\
    2, 3 & GUI Front End $\rightarrow$ Window and Platform Layer, Font Subsystem & The GUI draws text in native windows \\
    4 & Multiplexer $\rightarrow$ Terminal Emulation Core & Pane output updates the terminal model \\
    5 & Multiplexer $\rightarrow$ PTY & The local domain runs programs in PTYs \\
    6, 7 & Multiplexer $\rightarrow$ Mux Client, SSH Client & Remote and SSH domains supply panes \\
    8 & Mux Client $\rightarrow$ Multiplexer & The client registers remote panes \\
    9 & Mux Client $\rightarrow$ Mux Server & The server owns the remote sessions \\
    10, 11 & Mux Client, Mux Server $\rightarrow$ Mux Protocol & Both use the same message format \\
    12 & Mux Server $\rightarrow$ Multiplexer & The server applies client requests \\
    13 & Mux Client $\rightarrow$ SSH Client & SSH multiplexing domains use SSH \\
    14 & GUI Front End $\rightarrow$ Mux Server & The GUI serves \texttt{wezterm cli} \\
    15, 16 & GUI Front End, Multiplexer $\rightarrow$ Configuration and Lua Engine & Settings and Lua event handlers \\
    17, 18 & Configuration and Lua Engine $\rightarrow$ Multiplexer, GUI Front End & Scripts call their Lua APIs \\
    19--23 & Font Subsystem, Terminal Emulation Core, Mux Client, Mux Server, SSH Client $\rightarrow$ Configuration and Lua Engine & Each reads its settings, such as the colour palette \\
    \bottomrule
  \end{tabularx}
\end{table}
